\chapter{Tái lập và kiểm toán bằng chứng}\label{app:evidence}
\section{Baseline, nguồn và sổ bằng chứng}
Nguyên mẫu nền được nhận diện bằng phiên bản \BaselineCommit. Mỗi mã E01--E17 liên kết mệnh đề với ngày, phạm vi và giới hạn. Hồ sơ kiểm chứng đi kèm lưu đường dẫn mã, tên test, lệnh thực thi và checksum; phụ lục này trình bày nội dung cần để đánh giá kết luận. Dữ liệu đo thô và dữ liệu người tham gia được quản lý riêng.
\begin{longtable}{@{}P{.11\textwidth}P{.79\textwidth}@{}}
\caption{Sổ bằng chứng và giới hạn kiểm toán}\label{tab:evidence-register}\\
\toprule Mã & Mệnh đề, nguồn và provenance \\
\midrule\endfirsthead
\toprule Mã & Mệnh đề, nguồn và provenance \\
\midrule\endhead
\phantomsection\label{ev:E01}E01 & \textbf{Baseline và nguồn yêu cầu}. TÀI LIỆU VÀ MÃ. \newline Ngày: 2026-10-09. Phiên bản: e829f73. \newline Nguồn: Thuyết minh, đặc tả yêu cầu và ma trận quyền. \newline Kết quả: Đối chiếu yêu cầu và phạm vi. Thuyết minh ghi hai khoảng thời gian khác nhau.\newline Skip: Không áp dụng. \newline Giới hạn: Yêu cầu\slash phạm vi từ hồ sơ; thời hạn và phần mở rộng học thuật chưa được xác nhận. Không coi status runtime cũ là hiện tại. \\
\phantomsection\label{ev:E02}E02 & \textbf{Protocol và quản trị nghiên cứu}. TÀI LIỆU VÀ MÃ. \newline Ngày: 2026-08-25; đọc lại 2026-10-09. Phiên bản: CHƯA XÁC LẬP. \newline Nguồn: Protocol 0.1 và nhật ký tìm kiếm tài liệu. \newline Kết quả: Protocol 0.1 giữ RQ1–RQ4 và baseline Pool\slash Silo; note yêu cầu phụ lục cho phép đo ba placement.\newline Skip: Không áp dụng. \newline Giới hạn: Chưa khóa tìm kiếm và chưa có dữ liệu người dùng; báo cáo không sửa hồi tố protocol. \\
\phantomsection\label{ev:E03}E03 & \textbf{Artifact Surefire có sẵn trước lượt viết}. TỆP KẾT QUẢ. \newline Ngày: Quan sát 2026-10-09; mtime 2026-09-24 Asia\slash Taipei. Phiên bản: CHƯA XÁC LẬP. \newline Nguồn: Bộ kết quả Surefire được lưu trước lần kiểm tra mới. \newline Kết quả: Artifact cũ 106 test, 0 failure\slash error, 47 skip: 59 thực thi. Đã lưu bản sao trước kiểm tra mới.\newline Skip: 47: Docker không khả dụng. \newline Giới hạn: Mtime không chứng minh ngày chạy hoặc commit; Kết quả hiện tại thuộc lượt E04. \\
\phantomsection\label{ev:E04}E04 & \textbf{Kiểm tra backend mới và môi trường hiện tại}. KIỂM TRA MỚI. \newline Ngày: 2026-10-09 11:13:19 +08:00. Phiên bản: e829f73; không thay mã ứng dụng. \newline Nguồn: Biên bản backend ngày 09/10/2026 và bảng tổng hợp suite có checksum. \newline Kết quả: Fresh compile 121 source + 27 test source; test goal đạt: 106 tổng, 59 thực thi, 47 skip, 0 failure\slash error. Maven\slash test runtime Java 21.0.12.1; compiler javac 25.0.4.1 fork với --release 21.\newline Skip: 47: \Code{disabledWithoutDocker}. \newline Giới hạn: Verify chưa đạt do thiếu Maven plugin offline. Compiler trong runtime 21 báo release 21 unsupported, nên dùng compiler fork đã ghi rõ. Frontend\slash E2E không chạy vì thiếu Node\slash runtime stack. \\
\phantomsection\label{ev:E05}E05 & \textbf{Checkpoint EXT ba placement}. BIÊN BẢN LỊCH SỬ. \newline Ngày: 2026-09-08 UTC. Phiên bản: Working tree sau c226676; không gán HEAD hiện tại. \newline Nguồn: Biên bản mở rộng ba placement ngày 08/09/2026. \newline Kết quả: Application V8, backend 96\slash 96 (skip 0), frontend 16, Playwright 3 theo biên bản; schema isolation\slash upgrade local.\newline Skip: Backend 0 theo biên bản. \newline Giới hạn: Kết quả lịch sử local; extension không tự là phạm vi học thuật đã duyệt hoặc runtime hôm nay. \\
\phantomsection\label{ev:E06}E06 & \textbf{Checkpoint nhắc deadline V12}. BIÊN BẢN LỊCH SỬ. \newline Ngày: Nội dung 2026-09-19 UTC+7; tên file 2026-09-18. Phiên bản: an\_upgrade\_features theo biên bản; SHA chưa ghi. \newline Nguồn: Biên bản nhắc hạn có nội dung ngày 19/09/2026. \newline Kết quả: Biên bản ghi deadline integration 6\slash 6, backend 106, frontend 22, E2E 3, application V12; kiểm tra stale và exact assignee.\newline Skip: Tổng skip không ghi riêng. \newline Giới hạn: Phân biệt ngày nội dung\slash tên file; không suy 106\slash 106 pass hiện tại, không có SMTP Internet\slash VAPID thật. \\
\phantomsection\label{ev:E07}E07 & \textbf{P0 baseline local}. BIÊN BẢN LỊCH SỬ. \newline Ngày: 2026-08-26. Phiên bản: Theo biên bản; không liên kết HEAD hiện tại. \newline Nguồn: Biên bản xác minh baseline ngày 26/08/2026. \newline Kết quả: Baseline Pool\slash Silo, migration\slash role\slash RLS, API\slash Compose và k6 smoke theo hồ sơ.\newline Skip: Theo từng lệnh trong biên bản. \newline Giới hạn: Smoke timing không phải số đo nghiên cứu; trạng thái stack khi chốt không được giả định còn nguyên. \\
\phantomsection\label{ev:E08}E08 & \textbf{P1 provisioning và storage recovery}. BIÊN BẢN LỊCH SỬ. \newline Ngày: 2026-08-31. Phiên bản: Theo từng biên bản, không gán một SHA chung. \newline Nguồn: Các biên bản cấp phát và phục hồi lưu trữ ngày 31/08/2026. \newline Kết quả: Hồ sơ ghi MinIO dead letter attempt 5\slash requeue riêng Pool–Silo, force-kill JVM sau các ranh giới, lease recovery, rollback lỗi lặp và manual retry.\newline Skip: Theo từng lệnh trong biên bản. \newline Giới hạn: Không chạy lại fault injection trong lượt viết; không đại diện availability\slash durability production. \\
\phantomsection\label{ev:E09}E09 & \textbf{Sàng lọc guard-omission isolation}. SÀNG LỌC LỊCH SỬ. \newline Ngày: 2026-08-31. Phiên bản: Nền harness 2b430b7 theo ADR; quan sát lưu trong biên bản lượt 3. \newline Nguồn: Biên bản sàng lọc isolation ngày 31/08/2026 và ADR-0003. \newline Kết quả: Guarded contract 6\slash 6; 18 observation: 6 explicit\slash Hibernate Pool leak, 3 RLS Pool protected, 9 Silo boundary protected.\newline Skip: Theo hồ sơ; không chạy harness mới. \newline Giới hạn: Harness pass là phân loại đúng, không đổi leak thành PASS. Measurement\slash scorecard\slash hồ sơ loại checksum-backed chưa hoàn tất. \\
\phantomsection\label{ev:E10}E10 & \textbf{N1 tenant context}. MÃ VÀ THIẾT KẾ KIỂM TRA. \newline Ngày: 2026-10-09. Phiên bản: e829f73. \newline Nguồn: Các thành phần phân giải host, kiểm tra phiên và quản lý context. \newline Kết quả: Host\slash token\slash status\slash membership version binding; ThreadLocal finally cleanup; transaction-local set\_config và search\_path. Unit test hiện có chạy trong E04.\newline Skip: Xem E04. \newline Giới hạn: Source mô tả cơ chế; chưa có trusted proxy\slash DNS\slash TLS Internet được nghiệm thu. \\
\phantomsection\label{ev:E11}E11 & \textbf{N2 datasource và isolation}. MÃ VÀ THIẾT KẾ KIỂM TRA. \newline Ngày: 2026-10-09. Phiên bản: e829f73. \newline Nguồn: Các thành phần phân giải datasource, cấp phát và migration. \newline Kết quả: Ba placement; runtime role\slash schema\slash database, migration V1–V12 và guard\slash RLS.\newline Skip: RLS 4 + schema 1 skip trong E04. \newline Giới hạn: Runtime privilege cần kiểm tra trên database thật; source không chứng minh mọi deployment đúng. \\
\phantomsection\label{ev:E12}E12 & \textbf{N3 identity và authorization}. MÃ VÀ THIẾT KẾ KIỂM TRA. \newline Ngày: 2026-10-09. Phiên bản: e829f73. \newline Nguồn: Dịch vụ đăng nhập, thành viên và chính sách quyền project. \newline Kết quả: Account chung, tenant-bound session, active membership\slash security version và project role; deny\slash IDOR assertions hiện có.\newline Skip: 17 project integration skip; management unit thực thi. \newline Giới hạn: Chỉ kết luận các case có evidence; không chứng nhận identity production hoặc mọi ô policy. \\
\phantomsection\label{ev:E13}E13 & \textbf{N4 version, outbox, reminder và delivery}. MÃ VÀ THIẾT KẾ KIỂM TRA. \newline Ngày: 2026-10-09. Phiên bản: e829f73. \newline Nguồn: Dịch vụ mutation, outbox, nhắc hạn và hàng đợi delivery. \newline Kết quả: Version predicate, batch transaction\slash outbox, sổ nhắc hạn, stale recheck và email queue riêng; outbox duyệt tenant tuần tự.\newline Skip: 6 deadline integration skip trong E04. \newline Giới hạn: Không exactly-once SMTP; không kế thừa SKIP LOCKED provisioning cho outbox; multi-worker còn thiếu. \\
\phantomsection\label{ev:E14}E14 & \textbf{N5 provisioning và tài nguyên}. MÃ VÀ THIẾT KẾ KIỂM TRA. \newline Ngày: 2026-10-09. Phiên bản: e829f73. \newline Nguồn: Bộ claim/coordinator, provisioner, storage và giới hạn tài nguyên. \newline Kết quả: Atomic claim\slash lease, compensation, datasource eviction, process-local quota lock và tenant bucket limiter.\newline Skip: Integration Docker xem E04. \newline Giới hạn: Chưa đo global connection budget, quota\slash limiter multi-instance hoặc provider thật; pool cap logic chưa là proof hard bound toàn hệ thống. \\
\phantomsection\label{ev:E15}E15 & \textbf{Contract, architecture và workload}. TÀI LIỆU VÀ MÃ. \newline Ngày: 2026-10-09. Phiên bản: e829f73. \newline Nguồn: Hợp đồng OpenAPI, tài liệu kiến trúc và workload k6. \newline Kết quả: API chung; baseline constant-vus + think time; load constant-arrival-rate. ADR thiết kế concurrency khác worker tuần tự.\newline Skip: Không áp dụng. \newline Giới hạn: Workload đọc chưa là mixed Kanban; protocol ba placement và capability cần đăng ký bổ sung. \\
\phantomsection\label{ev:E16}E16 & \textbf{Frontend và E2E lịch sử}. LỊCH SỬ VÀ MÔI TRƯỜNG. \newline Ngày: 08\slash 09 và 19\slash 09; môi trường đọc lại 09\slash 10\slash 2026. Phiên bản: Theo E05\slash E06. \newline Nguồn: Biên bản frontend/E2E tại hai checkpoint E05 và E06. \newline Kết quả: 16 frontend test ở EXT, 22 ở mốc sau; 3 E2E theo hồ sơ. Node Linux không có trong môi trường hiện tại.\newline Skip: Không có lượt frontend\slash E2E mới. \newline Giới hạn: Không tuyên bố contract\slash lint\slash test\slash build hoặc E2E đã chạy lại hôm nay. \\
\phantomsection\label{ev:E17}E17 & \textbf{Xuất bản LaTeX\slash PDF}. XUẤT BẢN. \newline Ngày: 2026-10-09. Phiên bản: e829f73; nội dung báo cáo cập nhật riêng. \newline Nguồn: Biên bản xuất bản và kiểm tra bố cục LaTeX/PDF. \newline Kết quả: Doctor, all smoke và check đạt; review trực quan được ghi trong biên bản xuất bản.\newline Skip: Không áp dụng. \newline Giới hạn: Định dạng\slash build không thay cổng nghiên cứu, đủ trang nội dung hoặc nghiệm thu. Kết quả cuối đọc tại biên bản xuất bản. \\
\bottomrule\end{longtable}


\section{Điều kiện tái lập xuất bản}
Nguồn báo cáo sử dụng XeLaTeX, Biber hoặc BibTeX, Poppler và Times New Roman đủ bốn kiểu. Quy trình gồm kiểm tra công cụ, biên dịch đến khi tham chiếu ổn định, kiểm tra PDF và xem trực quan. Báo cáo một mặt, hai mặt, bản tin và tóm tắt Việt--Anh được xuất từ cùng nguồn; hướng dẫn lệnh và vị trí artifact nằm trong hồ sơ xuất bản đi kèm.

\section{Tái lập kiểm tra kỹ thuật và điều kiện}
Backend verification chuẩn dùng Java 21 và Maven wrapper với mục tiêu verify. Frontend cần Node ít nhất 22.12, dependency đúng lockfile, rồi contract\slash lint\slash test\slash build. Testcontainers cần Docker; E2E cần stack và environment đúng. Không giả định môi trường hiện tại có các điều kiện ấy chỉ vì một checkpoint trước đây đã pass.

Lượt viết này dùng cache Maven sao chép trong thư mục tạm để tránh ghi home read-only. Các lệnh và kết quả thực tế, gồm mục tiêu nào thất bại và test nào skip, được lưu ở biên bản E04. Để đóng baseline cuối cần chạy lại verify và frontend đầy đủ trên môi trường chuẩn, kèm commit, ngày, số test và skipped count.

\section{Manifest cho đợt thực nghiệm sau}
Một manifest tối thiểu ghi run ID\slash class, protocol ID, commit\slash image digest, working-tree state, seed, tenant\slash account\slash project\slash task counts, placement, VU hoặc arrival rate, think time, warm-up\slash cửa sổ đo, pool\slash limiter và phiên bản công cụ. Phần cứng, dịch vụ chạy chung, metric nguồn và checksum phải đi cùng raw artifact. Thông tin token\slash secret không nằm trong manifest.

Kết quả derived được sinh từ raw input đã QA. Một run development\slash pilot không đổi nhãn thành experiment; mọi điều chỉnh protocol phải ghi trước. Các bảng pending trong Chương 4 không được điền tay từ ảnh dashboard hoặc từ một smoke iteration.

\section{Checklist xuất bản và nội dung}
Checklist gồm đủ Mở đầu\slash bốn chương\slash Kết luận\slash bibliography\slash phụ lục; N1--N5 có đủ chuỗi sáu bước; mỗi RQ có bằng chứng\slash giới hạn\slash công việc bổ sung; citation\slash reference ổn định; bảng dài lặp header; năm sơ đồ đọc được; Times New Roman nhúng và A4. Kiểm tra tự động không thay xem trực quan bìa, mục lục, hình và bibliography. Số trang chính và mức đáp ứng bản cuối được ghi riêng trong biên bản xuất bản.
